% !TeX spellcheck = ru_RU
% !TEX root = report.tex

\section{Реализация}
\label{sec:impl}

Программная часть состоит из реализации на Python (\code{hypergraph\_isomorphism.py}, 634 строки), реализации на \Cpp{} (\code{hypergraph\_isomorphism.cpp}, 993 строки), сценария сравнения \code{compare.py} (561 строка) и файла \code{examples.txt} с примерами статьи.
Обе реализации читают задачи в одном текстовом формате: строки \code{pred p2 2 n} объявляют предикаты (арность, тип рёбер), строки \code{g1 p2 a1 a2} и \code{g2 ...}~--- рёбра $G_1$ и $G_2$.
Шаблон $H$ получается из $G_2$ заменой вершин, упорядоченных \enquote{естественно} ($b_2 < b_{10}$), на $x_1, x_2, \dots$; ищутся все унификаторы $\lambda$ с $H\lambda = G_1$, и каждый печатается как изоморфизм $G_2 \to G_1$.
Части программы названы так же, как в статье: характеристики, структурное отображение, \ISOMm{}, \ISOMone{}, \proc{MAP}, \proc{BESTMAP}, \proc{CONSISTENCY}, \proc{REDUCTION} и цикл \ISOMn{}; для проверки добавлен полный перебор.

\subsection{Реализация на Python}
\label{subsec:python}

\paragraph{Данные.}
Предикат~--- неизменяемый класс \code{Pred} (имя, арность, коммутативность), ребро \code{Edge}~--- предикат и кортеж (тип~(1)) или \code{frozenset} (тип~(2)) аргументов, гиперграф~--- \code{Counter} рёбер.
Переменные шаблона~--- объекты класса \code{Var}, подстановка~--- словарь \enquote{переменная $\to$ константа}.

\paragraph{Структурное отображение.}
Функция \code{build\_mapping} строит отображение между $G^i$ и $H^i\lambda$, где часть переменных уже заменена константами, только если выполнены необходимые условия: число рёбер и мультимножества характеристик совпадают, а каждая константа, уже стоящая в $H^i\lambda$, имеет в $G^i$ ту же характеристику.
Иначе ветвь поиска отсекается.

\paragraph{\ISOMm{}.}
Подстановки перебираются по классам: для класса $k$ перебираются перестановки его констант, и для каждой рекурсивно обрабатываются следующие классы (листинг~\ref{lst:isom1}).
Ветвь отсекается, как только ребро $H^i\lambda$ без переменных не встречается в $G^i$ (аналог шагов~8--12 \ISOMone{}), причём каждое ребро проверяется один раз~--- после назначения того класса, на котором в нём не остаётся переменных.
На листе дополнительно проверяется $H^i\lambda = G^i$ (шаг~7).

\begin{lemma}
    \label{lem:levels}
    Проверка по уровням отсекает те же подстановки, что и проверка всех рёбер без переменных после каждого класса.
\end{lemma}
\begin{proof}
    Ребро, ставшее рёбром без переменных на классе $k' < k$, уже проверено с теми же значениями переменных.
    Константы, уже стоящие в $H^i\lambda$, исключены из классов, а внутри класса значения образуют перестановку, поэтому подстановка инъективна на вершинах $H^i\lambda$ и разные рёбра переходят в разные; значит, кратности рёбер можно сравнивать поштучно.
\end{proof}

\begin{lstlisting}[style=code, language=Python, caption={Ядро \ISOMm{} на Python}, label={lst:isom1}]
level = {x: k + 1 for k, (xs, _) in enumerate(classes) for x in xs}
checks = [[] for _ in range(len(classes) + 1)]
for e, m in H_i.items():   # уровень ребра: класс, где кончаются переменные
    k = max((level[v] for v in e.args if is_var(v)), default=0)
    checks[k].append((e.args, e.pred.commutative, m))

def rec(k):
    if k == len(classes):
        if substitute_edges(H_i, s) == G_i:      # шаг 7
            result.append(dict(s))
        return
    xs, cs = classes[k]
    for perm in itertools.permutations(cs):
        s.update(zip(xs, perm))
        if fits(checks[k + 1]):
            rec(k + 1)
    for x in xs:
        del s[x]
\end{lstlisting}

\paragraph{\ISOMn{}.}
Очередь~--- двоичная куча \code{heapq} с ключом \enquote{число возможных подстановок, номер добавления}.
\proc{MAP} и \proc{BESTMAP} вызываются при добавлении вершины в очередь, так как её приоритет~--- число возможных подстановок $\prod t_k!$ выбранного отображения; при равенстве выбирается предикат с меньшим именем.
\proc{REDUCTION} подставляет унификатор только в подграфы, все переменные которых уже заданы.
С параметром \code{verbose=True} печатается журнал поиска, по которому можно восстановить дерево (как на рис.~\ref{fig:tree}).

\begin{lemma}
    \label{lem:tree}
    Если \proc{BESTMAP} выбирает подграф однозначно по $\lambda$, то все вершины дерева \ISOMn{} имеют различные частичные унификаторы, а все найденные унификаторы различны.
\end{lemma}
\begin{proof}
    Дети вершины с унификатором $\lambda$ имеют вид $\lambda\lambda_t$, где все $\lambda_t$ определены на одном множестве $V$ переменных $H^{i_0}\lambda$ и попарно различны.
    У потомков унификатор только расширяется, поэтому вершины из поддеревьев разных детей различаются на $V$, а вершина и её потомок~--- областью определения.
    Полные унификаторы получаются из пар \enquote{вершина, кандидат} и в очередь не попадают, поэтому они тоже различны.
\end{proof}

Поэтому реализация не хранит множество пройденных состояний, а отсутствие повторов среди унификаторов проверяется при сравнении.
Самопроверка модуля сравнивает \ISOMn{} и \ISOMone{} с полным перебором на 400 случайных парах гиперграфов.

\subsection{Реализация на \Cpp{}}
\label{subsec:cpp}

Реализация на \Cpp{} повторяет Python-версию по шагам, но использует компактные структуры данных.
Константы и переменные нумеруются целыми числами в естественном порядке имён, подстановка~--- массив, где $-1$ означает \enquote{не задано}.
Рёбра хранятся как строки длины $k$ (у коммутативных строка отсортирована), рёбра $G^i$~--- в хеш-таблице с открытой адресацией, характеристики~--- строки длины $k$ или~1, сравниваемые лексикографически, как кортежи в Python.
Число подстановок $\prod t_k!$ считается точно в 128-битном целом и при переполнении насыщается.
Гиперграф считается множеством рёбер, как в статье.
Чтобы деревья поиска совпадали, совпадают все порядки, от которых зависит обход (табл.~\ref{tab:order}).

\begin{table}[ht]
    \centering
    \small
    \caption{Порядки, от которых зависит обход, и их реализация}
    \label{tab:order}
    \begin{tabularx}{\textwidth}{@{}>{\raggedright\arraybackslash}p{3.6cm} >{\raggedright\arraybackslash}X >{\raggedright\arraybackslash}X@{}}
        \toprule
        Что упорядочено                     & Python                                                     & \Cpp{}                                                         \\
        \midrule
        перестановки класса                 & \code{itertools.permutations} от отсортированного списка   & \code{std::next\_permutation} от возрастающего массива         \\
        очередь                             & \code{heapq}, ключ ($\prod t_k!$, счётчик)                    & \code{std::push\_heap}, тот же ключ                            \\
        предикаты в \proc{BESTMAP} и \proc{REDUCTION} & по строке вида \code{p2\^{}2}                     & по той же строке (байты UTF-8 упорядочены как коды символов)   \\
        классы отображения                  & по характеристике, затем по имени                          & по строке характеристики, затем по номеру                      \\
        \bottomrule
    \end{tabularx}
\end{table}

Программа печатает для каждой задачи ответ, время (все замеры серии, их среднее и стандартное отклонение), статистику поиска и найденные изоморфизмы; опция \code{-v} печатает журнал в том же формате, что и Python-версия, \code{--algo} выбирает \ISOMn{}, \ISOMone{} или полный перебор.
Сборка: \code{g++ -O2 -std=c++17}.

\subsection{Средства сравнения}
\label{subsec:compare}

Сценарий \code{compare.py} генерирует наборы задач (раздел~\ref{subsec:families}), запускает обе реализации и для каждой задачи проверяет, что:
\begin{itemize}
    \item совпадают ответ и множества найденных изоморфизмов, а повторов среди унификаторов нет;
    \item совпадает статистика дерева поиска: вызовы \proc{MAP} и отсечения на нём, число вершин (вызовов \ISOMm{}), суммы $r$ и $T$, число полных и частичных унификаторов; в Python она собирается обёртками, временно подставленными вместо функций модуля;
    \item на малых задачах полный перебор обеих реализаций и \ISOMone{} (если предикат один) дают то же множество, что и \ISOMn{};
    \item на примерах статьи журналы поиска совпадают построчно.
\end{itemize}
Время замеряется для поиска всех унификаторов \ISOMn{} без чтения входа, по методике из раздела~\ref{subsec:conditions}; сырые замеры сохраняются для анализа.
